% Chapter 1

\chapter{Introducción general} % Main chapter title

\label{Chapter1} % For referencing the chapter elsewhere, use \ref{Chapter1} 
\label{IntroGeneral}

En este capítulo se presentan los fundamentos de las tecnologías de interfaz cerebro-computador y se describe el estado actual de la investigación en el área. Asimismo, se exponen la motivación, los objetivos y el alcance definidos para el desarrollo del proyecto.

%----------------------------------------------------------------------------------------

% Define some commands to keep the formatting separated from the content 
\newcommand{\keyword}[1]{\textbf{#1}}
\newcommand{\tabhead}[1]{\textbf{#1}}
\newcommand{\code}[1]{\texttt{#1}}
\newcommand{\file}[1]{\texttt{\bfseries#1}}
\newcommand{\option}[1]{\texttt{\itshape#1}}
\newcommand{\grados}{$^{\circ}$}

%----------------------------------------------------------------------------------------

%\section{Introducción}

%----------------------------------------------------------------------------------------
\section{Interfaces cerebro-computador}

Las interfaces cerebro-computador o \textit{BCI} (del inglés, \textit{Brain-Computer Interface}) se definen como sistemas que establecen una vía de comunicación directa entre el cerebro humano y un dispositivo externo \cite{wolpaw2002bci}. Estos sistemas permitieron la traducción de patrones de actividad cerebral en comandos de control sin la intervención de las vías periféricas neuromusculares tradicionales.

El funcionamiento de una \textit{BCI} se fundamentó en la adquisición de señales bioeléctricas, su procesamiento digital y la posterior ejecución de acciones en un entorno computacional o físico. En el desarrollo de este proyecto, se consideraron las señales electroencefalográficas o \textit{EEG} (del inglés, \textit{electroencephalogram}) debido a su naturaleza no invasiva y su capacidad para registrar la intención motriz del usuario en tiempo real. La integración de estos dispositivos con servicios en la nube facilita el procesamiento de grandes volúmenes de datos, optimizando la capacidad de respuesta del los sistemas ante tareas de imaginación motriz.


\subsection{Arquitectura funcional de un sistema BCI}

El modelo operativo de una \textit{BCI} inicia con una cadena de procesamiento compuesta por etapas secuenciales. En la fase inicial, se realiza la adquisición de señales mediante sensores electroencefalográficos colocados sobre el cuero cabelludo. Luego, los datos crudos se somenten a procesos de pre-procesamiento para la eliminación de artefactos, tales como parpadeos o movimientos musculares, mediante el uso de filtros digitales.
Posteriormente, se realiza la extracción de características que permite identificar los componentes relevantes de la señal asociados a la intención del usuario. Finalmente, se implementan algoritmos de clasificación para traducir dichas características en comandos de control.


\subsubsection{Imaginación motriz y ritmos sensorimotores}

La imaginación motriz se define como el proceso cognitivo en el cual un individuo simula mentalmente una acción motora sin ejecutar un movimiento físico real \cite{pfurtscheller2001motor}. Este fenómeno produce cambios detectables en los ritmos sensoriales de la corteza cerebral, específicamente en las bandas de frecuencia mu (8-13 Hz) y beta (13-30 Hz).

\subsection{Dispositivos de adquisición de señales EEG}

La captura de la actividad eléctrica cerebral se realiza habitualmente mediante sistemas de \textit{EEG}, los cuales se clasifican según su movilidad y el tipo de sensores empleados. En el ámbito de la investigación, se destacan dispositivos de grado médico y equipos portátiles diseñados para aplicaciones fuera del entorno clínico \cite{casson2010wearable}.

Aunque en este proyecto se utilizó un conjunto de datos o dataset previamente registrado, la comprensión de estos dispositivos resultó pertinente para el análisis de la calidad de la señal.

%----------------------------------------------------------------------------------------

\section{Estado del arte}

La presente revisión del estado del arte tiene como objetivo identificar las metodologías dominantes en la clasificación de tareas de imaginación motriz mediante \textit{EEG}. En los últimos años, el estudio de las BCI ha cobrado relevancia debido a la convergencia entre la neurociencia computacional y el aprendizaje profundo, lo cual permite aprender directamente de las señales crudas,
sin depender de métodos tradicionales de extracción de características como el patrón especial común o \textit{CSP} (del inglés, \textit{Common Spatial Pattern})\cite{schirrmeister2017deep}.

Las redes neuronales convolucionales o \textit{CNN} (del ingles, \textit{Convolutional Neural Network}) constituyen una de las aproximaciones más utilizadas. Estas redes capturan patrones espaciales en las señales \textit{EEG} y han demostrado un buen desempeño en la clasificación de tareas de imaginación motriz. Modelos como EEGNet introducen arquitecturas compactas que logran buenos resultados incluso con pocos datos de entrenamiento \cite{lawhern2018eegnet}.

Por otra parte, las redes recurrentes, especialmente las de memoria a corto y largo plazo o \textit{LSTM} (del inglés, \textit{Long Short-Term Memory}), permiten modelar la dinámica temporal de las señales \textit{EEG}. Estas arquitecturas capturan dependencias en el tiempo que resultan clave para identificar patrones asociados a la imaginación de movimiento. Estudios recientes reportan que modelos basados en LSTM alcanzan desempeños superiores frente a métodos tradicionales \cite{roy2019deep}.

Asimismo, las arquitecturas híbridas que combinan \textit{CNN} y \textit{LSTM} han ganado relevancia. Estas combinaciones permiten integrar información espacial y temporal en un mismo modelo. Investigaciones muestran que este enfoque mejora la capacidad de generalización y aumenta la precisión en comparación con modelos individuales  \cite{bashivan2015learning,roy2019deep}.

En los trabajos más recientes, se han propuesto variantes más eficientes como redes convolucionales ligeras y modelos que operan en el dominio frecuencial. Estas propuestas buscan reducir la complejidad del modelo y mejorar el rendimiento con conjuntos de datos limitados, lo cual representa un desafío común en aplicaciones de interfaces cerebro-computador \cite{craik2019deep}.

La literatura técnica muestra un interés  en el desplazamiento desde los métodos basados en la extracción manual de características, como los \textit{CSP}, hacia sistemas de aprendizaje profundo. Mientras que los algoritmos tradicionales destacan por su baja carga computacional, las investigaciones recientes exponen que las arquitecturas híbridas logran una mayor robustez ante la variabilidad entre sujetos. Sin embargo, aún persisten
retos relacionados con la variabilidad entre usuarios, la necesidad de grandes volúmenes de datos y la robustez en entornos reales \cite{roy2019deep}.

A continuación, se presenta una tabla comparativa técnica de las arquitecturas y enfoques identificados en la revisión bibliográfica

\begin{table}[h]
	\centering
	\caption[]{Comparativa de enfoques tecnológicos para la clasificación de imaginación motriz.}
	\begin{tabular}{p{2cm} p{2cm} p{3.5cm} p{3.5cm}}
		\toprule
		\textbf{Enfoque} & \textbf{Técnicas} & \textbf{Ventajas} & \textbf{Limitaciones} \\
		\midrule
		Estadístico           & CSP + SVM    & Alta interpretabilidad y eficiencia en sistemas embebidos.  & Sensible a artefactos y requiere calibración manual. \\\\
		Redes Convolucionales & CNN (EEGNet) & Reducción de parámetros y extracción automática de rasgos. & Necesidad de grandes volúmenes de datos etiquetados. \\\\
		Arquitectura Híbrida  & CNN + LSTM   & Modelado de dependencias temporales y espaciales. & Complejidad computacional elevada y riesgo de sobreajuste. \\\\
		\bottomrule
	\end{tabular}
	\label{tab:peces}
\end{table}
%----------------------------------------------------------------------------------------
\newpage
\section{Motivación}

El desarrollo de este proyecto se justificó en la necesidad de mejorar la precisión en la predicción de la intención de movimiento de un usuario a partir de señales \textit{EEG} asociadas a tareas de imaginación motriz. A pesar de los avances en el procesamiento de señales biológicas, los métodos tradicionales presentaron dificultades para adaptarse a la variabilidad de los ritmos sensorimotores entre diferentes individuos, lo que limitó su aplicación en entornos reales fuera del laboratorio \cite{wolpaw2002bci,lotte2015signal}.

La implementación de modelos de inteligencia artificial avanzados permitió abordar la naturaleza no lineal y compleja de las señales de imaginación motriz, las cuales presentan una alta variabilidad entre sujetos. Durante el análisis preliminar, se identificó que el procesamiento local en dispositivos embebidos restringía el uso de arquitecturas de aprendizaje profundo debido a las limitaciones de memoria y capacidad de cómputo de los sistemas portátiles. Por lo tanto, la integración de una infraestructura en la nube se presentó como una solución técnica viable para ejecutar modelos de alta complejidad sin comprometer la portabilidad ni la autonomía del sistema de adquisición.

Asimismo, la motivación de este trabajo residió en la creación de una arquitectura escalable que facilitara el entrenamiento continuo de los modelos. Al centralizar el procesamiento, se permitió la acumulación de datos para el refinamiento de los algoritmos, buscando una reducción significativa en el tiempo de calibración requerido para cada usuario. Esta aproximación tecnológica no solo optimizó el rendimiento predictivo, sino que también sentó las bases para interfaces cerebro-computador más robustas y eficientes en la práctica clínica y asistencial.


\newpage
%----------------------------------------------------------------------------------------

\section{Objetivos y alcance}

El desarrollo de este proyecto se orientó a la creación de una infraestructura tecnológica capaz de procesar y clasificar señales \textit{EEG} en un entorno distribuido. A continuación, se presentan el objetivo general, los objetivos especificos y el alcance del proyecto.

\subsection{Objetivo general}

Establecer una base técnica que permitiera analizar patrones neuronales asociados a tareas de imaginación motriz, integrando modelos de inteligencia artificial y servicios en la nube para servir como referencia en futuros desarrollos de interfaces cerebro-computador.

\subsection{Objetivos específicos}
\label{subsec:configurando}

\begin{itemize}
	\item Implementar una arquitectura de red neuronal capaz de identificar con precisión las intenciones de movimiento a partir de señales \textit{EEG}.
	\item Diseñar un flujo de datos que integrara la adquisición de señales con el procesamiento remoto en una plataforma en la nube.
	\item Evaluar el desempeño del sistema mediante métricas de precisión y tiempo de respuesta en la clasificación de actividades de imaginación motriz.
	\item Documentar los procedimientos de pre-procesamiento y entrenamiento para facilitar la replicabilidad del modelo en investigaciones posteriores.
\end{itemize}

\subsection{Alcance}

El alcance de este trabajo comprendió desde la selección y limpieza del conjunto de datos hasta el despliegue del modelo de inferencia en servidores remotos. Se limitó el estudio al análisis de tareas motoras binarias como el movimiento imaginado de mano izquierda y derecha, utilizando datasets de acceso público para garantizar la validación estadística de los resultados.


El sistema se diseñó para operar de forma asíncrona, priorizando la precisión del clasificador sobre la respuesta inmediata de baja latencia, estableciendo así un marco de trabajo robusto para la investigación de algoritmos complejos. No se incluyó en esta etapa el desarrollo de hardware de adquisición propio ni la implementación de biorretroalimentación en tiempo real para el usuario final.